\documentclass[12pt,a4paper]{article}

\usepackage[T2A]{fontenc}
\usepackage[utf8]{inputenc}
\usepackage[russian]{babel}
\usepackage{amsmath,amssymb}
\usepackage{enumitem}
\usepackage{geometry}
\usepackage{hyperref}

\geometry{margin=2cm}
\setlist{nosep}
\hypersetup{colorlinks=true, linkcolor=blue, urlcolor=blue}

\DeclareMathOperator{\End}{End}
\DeclareMathOperator{\Ker}{ker}
\DeclareMathOperator{\Image}{im}
\DeclareMathOperator{\rank}{rank}
\DeclareMathOperator{\GL}{GL}
\DeclareMathOperator{\Span}{span}

\newcommand{\F}{\mathbb{F}}
\newcommand{\R}{\mathbb{R}}
\renewcommand{\C}{\mathbb{C}}
\newcommand{\Q}{\mathbb{Q}}
\newcommand{\Z}{\mathbb{Z}}
\newcommand{\Lin}[1]{\left\langle #1 \right\rangle}
\newcommand{\ds}{\displaystyle}

\begin{document}

\begin{center}
    {\Large \textbf{Алгебра, 2 семестр}}\\[4pt]
    {\large Краткий LaTeX-конспект по экзаменационным билетам 1--37}
\end{center}

\tableofcontents
\newpage

\section*{Как пользоваться конспектом}
\addcontentsline{toc}{section}{Как пользоваться конспектом}

Это сокращенная версия под экзаменационные билеты. В каждом билете оставлены определения, формулировки теорем, основные доказательства и алгоритмы. Исторические вставки, длинные примеры, олимпиадные упражнения, минимальный многочлен, корневые подпространства, идемпотентные операторы, разложение Фиттинга и триангуляризуемые операторы можно не учить, если преподаватель отдельно не требует.

\section{Теорема о делении многочлена с остатком. Свойства делимости}

Пусть $A$ --- целостное кольцо, $f,g\in A[x]$, причем старший коэффициент $g$ обратим в $A$.

\subsection*{Теорема о делении с остатком}

Существуют и единственны многочлены $q,r\in A[x]$, такие что
\[
    f=qg+r,\qquad \deg r<\deg g.
\]

\subsection*{Идея доказательства}

Существование доказывается индукцией по $\deg f$. Если $\deg f<\deg g$, то $q=0$, $r=f$. Если $\deg f\ge \deg g$, вычитаем из $f$ подходящий старший член, кратный $g$, чтобы степень уменьшилась, и применяем индукцию.

Единственность: если
\[
    f=qg+r=q_1g+r_1,
\]
то
\[
    (q-q_1)g=r_1-r.
\]
Если $q\ne q_1$, то степень левой части не меньше $\deg g$, а степень правой меньше $\deg g$, противоречие. Значит $q=q_1$ и $r=r_1$.

\subsection*{Делимость}

В целостном кольце $K$ говорят, что $a$ делит $b$, если существует $c\in K$, такое что
\[
    b=ac.
\]
Обозначение: $a\mid b$.

\subsection*{Свойства делимости}

\begin{enumerate}
    \item Если $a\mid b$ и $b\mid c$, то $a\mid c$.
    \item Если $a\mid b$ и $a\mid c$, то $a\mid (b\pm c)$.
    \item Если $a\mid b$, то $a\mid bc$ для любого $c$.
\end{enumerate}

Два элемента называются \textbf{ассоциированными}, если $a\mid b$ и $b\mid a$. В кольце многочленов над полем это означает, что они отличаются на ненулевой постоянный множитель.

\section{НОД двух многочленов. Теорема о существовании}

Пусть $P$ --- поле, $f,g\in P[x]$.

\subsection*{Определение}

Многочлен $d$ называется \textbf{наибольшим общим делителем} $f$ и $g$, если:
\begin{enumerate}
    \item $d\mid f$ и $d\mid g$;
    \item если $c\mid f$ и $c\mid g$, то $c\mid d$.
\end{enumerate}

НОД определен с точностью до ненулевого постоянного множителя. Часто выбирают унитарный НОД.

\subsection*{Алгоритм Евклида}

Пусть $\deg f>\deg g$. Последовательно делим с остатком:
\[
    f=gq_1+r_1,\qquad \deg r_1<\deg g,
\]
\[
    g=r_1q_2+r_2,\qquad \deg r_2<\deg r_1,
\]
\[
    r_{k-2}=r_{k-1}q_k+r_k.
\]
Степени остатков строго убывают, поэтому процесс заканчивается. Последний ненулевой остаток $r_k$ является НОД:
\[
    (f,g)=r_k.
\]

\subsection*{Теорема о существовании НОД}

Для любых ненулевых $f,g\in P[x]$ существует НОД. Он находится алгоритмом Евклида.

\section{Представимость НОД в виде $fu+gv$}

\subsection*{Теорема Безу для НОД}

Пусть $f,g\in P[x]$, где $P$ --- поле, и $d=(f,g)$. Тогда существуют $u,v\in P[x]$, такие что
\[
    d=fu+gv.
\]

\subsection*{Доказательство}

В алгоритме Евклида последний ненулевой остаток $r_k$ равен НОД. Из последнего деления выражаем $r_k$ через два предыдущих остатка, затем каждый предыдущий остаток выражаем через еще более ранние. Подставляя назад, получаем
\[
    r_k=fu+gv.
\]
Так как $r_k=(f,g)$, то
\[
    (f,g)=fu+gv.
\]

\subsection*{Следствие}

Многочлены $f$ и $g$ взаимно просты тогда и только тогда, когда существуют $u,v\in P[x]$, такие что
\[
    fu+gv=1.
\]

\section{Взаимно простые многочлены}

\subsection*{Определение}

Многочлены $f,g\in P[x]$ называются \textbf{взаимно простыми}, если
\[
    (f,g)=1.
\]

\subsection*{Критерий взаимной простоты}

$f$ и $g$ взаимно просты тогда и только тогда, когда существуют $u,v\in P[x]$, такие что
\[
    fu+gv=1.
\]

\subsection*{Свойства}

Пусть $a,b,c\in P[x]$.

\begin{enumerate}
    \item Если $(a,b)=1$ и $(a,c)=1$, то $(a,bc)=1$.
    \item Если $a\mid bc$ и $(a,b)=1$, то $a\mid c$.
    \item Если $a\mid c$, $b\mid c$ и $(a,b)=1$, то $ab\mid c$.
\end{enumerate}

\subsection*{Доказательство свойства 2}

Из $(a,b)=1$ следует, что существуют $u,v$, такие что
\[
    au+bv=1.
\]
Умножим на $c$:
\[
    acu+bcv=c.
\]
Так как $a\mid ac$ и $a\mid bc$, то $a\mid c$.

\section{Неприводимые и приводимые многочлены}

\subsection*{Определение}

Ненулевой многочлен $f\in P[x]$ положительной степени называется \textbf{неприводимым над $P$}, если он не раскладывается в произведение двух многочленов положительной степени из $P[x]$.

Если такое разложение существует, то $f$ называется \textbf{приводимым}.

\subsection*{Свойства}

\begin{enumerate}
    \item Любой многочлен первой степени неприводим.
    \item Неприводимость зависит от поля. Например, $x^2+1$ неприводим над $\R$, но приводим над $\C$:
    \[
        x^2+1=(x-i)(x+i).
    \]
    \item Если $p$ неприводим и $p\mid fg$, то $p\mid f$ или $p\mid g$.
\end{enumerate}

\subsection*{Примеры}

\begin{itemize}
    \item Над $\C$ неприводимы только многочлены первой степени.
    \item Над $\R$ неприводимы многочлены первой степени и квадратные многочлены с отрицательным дискриминантом.
    \item Над $\Q$ существуют неприводимые многочлены любой степени, например $x^n+2$ неприводим по критерию Эйзенштейна.
\end{itemize}

\subsection*{Критерий Эйзенштейна}

Пусть
\[
    f(x)=x^n+a_1x^{n-1}+\cdots+a_n\in\Z[x].
\]
Если существует простое число $p$, такое что
\[
    p\mid a_1,\ldots,a_n,\qquad p^2\nmid a_n,
\]
то $f$ неприводим над $\Z$, а значит и над $\Q$.

\section{Основная теорема о разложении многочлена на множители}

\subsection*{Теорема}

Если $P$ --- поле, то кольцо $P[x]$ факториально: всякий ненулевой многочлен $f\in P[x]$ единственным образом раскладывается в произведение неприводимых многочленов:
\[
    f=a\,p_1p_2\cdots p_s,
\]
где $a\in P^\times$, а $p_i$ --- унитарные неприводимые многочлены. Единственность понимается с точностью до порядка множителей.

\subsection*{Идея доказательства}

Существование получается разложением приводимых множителей до тех пор, пока все множители не станут неприводимыми. Процесс конечен, так как степени уменьшаются.

Единственность основана на свойстве неприводимого многочлена:
\[
    p\mid fg\Rightarrow p\mid f\ \text{или}\ p\mid g.
\]
Сравнивая два разложения, сокращаем одинаковые неприводимые множители.

\section{Корни многочленов. Теорема Безу. Схема Горнера}

\subsection*{Корень многочлена}

Элемент $c$ называется \textbf{корнем} многочлена $f(x)$, если
\[
    f(c)=0.
\]

\subsection*{Теорема Безу}

Для $f\in P[x]$ элемент $c\in P$ является корнем $f$ тогда и только тогда, когда
\[
    x-c\mid f(x).
\]

\subsection*{Доказательство}

По теореме о делении с остатком:
\[
    f(x)=(x-c)q(x)+r,
\]
где $r$ --- константа. Подставим $x=c$:
\[
    f(c)=r.
\]
Значит $f(c)=0$ тогда и только тогда, когда $r=0$, то есть $x-c$ делит $f$.

\subsection*{Следствия}

\begin{enumerate}
    \item Если $c_1,\ldots,c_s$ --- различные корни $f$, то
    \[
        (x-c_1)\cdots(x-c_s)\mid f(x).
    \]
    \item Ненулевой многочлен степени $n$ над целостным кольцом имеет не больше $n$ корней.
    \item Если два многочлена степени не выше $n$ совпадают в $n+1$ различных точках, то они равны.
\end{enumerate}

\subsection*{Схема Горнера}

Пусть
\[
    f(x)=a_0x^n+a_1x^{n-1}+\cdots+a_n.
\]
При делении на $x-c$ коэффициенты частного находятся рекуррентно:
\[
    b_0=a_0,\qquad b_k=a_k+cb_{k-1}.
\]
Остаток равен
\[
    f(c)=a_n+cb_{n-1}.
\]

\section{Кратные корни и производная}

\subsection*{Определение}

Элемент $c$ называется корнем кратности $k$ многочлена $f$, если
\[
    (x-c)^k\mid f(x),\qquad (x-c)^{k+1}\nmid f(x).
\]
То есть
\[
    f(x)=(x-c)^kg(x),\qquad g(c)\ne0.
\]

\subsection*{Критерий кратного корня}

$c$ является кратным корнем $f$ тогда и только тогда, когда
\[
    f(c)=0,\qquad f'(c)=0.
\]

\subsection*{Теорема о понижении кратности}

Пусть $P$ --- поле характеристики $0$, а $p(x)$ --- неприводимый множитель кратности $k$ в $f(x)$:
\[
    f(x)=p(x)^kg(x),\qquad (p,g)=1.
\]
Тогда $p(x)$ входит в $f'(x)$ с кратностью $k-1$.

В частности, если $c$ --- корень $f$ кратности $k$, то $c$ является корнем $f'$ кратности $k-1$.

\subsection*{Следствие}

Если $P$ имеет характеристику $0$, то
\[
    (f,f')=p_1^{k_1-1}\cdots p_r^{k_r-1},
\]
если
\[
    f=a\,p_1^{k_1}\cdots p_r^{k_r}
\]
--- разложение на различные неприводимые множители.

\section{Основная теорема алгебры и следствия}

\subsection*{Алгебраически замкнутое поле}

Поле $P$ называется \textbf{алгебраически замкнутым}, если всякий многочлен положительной степени из $P[x]$ имеет корень в $P$. Эквивалентно: всякий многочлен из $P[x]$ раскладывается на линейные множители.

\subsection*{Основная теорема алгебры}

Поле комплексных чисел $\C$ алгебраически замкнуто.

То есть любой многочлен положительной степени с комплексными коэффициентами имеет хотя бы один комплексный корень.

\subsection*{Следствия}

\begin{enumerate}
    \item Любой многочлен степени $n$ над $\C$ раскладывается в произведение $n$ линейных множителей:
    \[
        f(x)=a(x-c_1)\cdots(x-c_n).
    \]
    \item Над $\C$ неприводимы только многочлены первой степени.
    \item Многочлен с комплексными коэффициентами степени $n$ имеет ровно $n$ корней с учетом кратностей.
\end{enumerate}

\section{Невещественные корни. Неприводимость над $\C$ и $\R$}

\subsection*{Теорема о сопряженных корнях}

Пусть $f\in\R[x]$ и $z=a+bi$ --- корень $f$. Тогда сопряженное число
\[
    \overline z=a-bi
\]
тоже является корнем $f$.

\subsection*{Доказательство}

Так как коэффициенты $f$ вещественные, то
\[
    f(\overline z)=\overline{f(z)}.
\]
Если $f(z)=0$, то
\[
    f(\overline z)=\overline 0=0.
\]

\subsection*{Разложение над $\R$}

Если $z=a+bi$, $b\ne0$, то паре корней $z,\overline z$ соответствует вещественный квадратный множитель
\[
    (x-z)(x-\overline z)=x^2-2ax+(a^2+b^2),
\]
у которого дискриминант отрицателен:
\[
    D=-4b^2<0.
\]

\subsection*{Неприводимые многочлены}

\begin{itemize}
    \item Над $\C$ неприводимы только многочлены первой степени.
    \item Над $\R$ неприводимы только многочлены первой степени и квадратные многочлены с отрицательным дискриминантом.
\end{itemize}

\section{Линейное пространство}

\subsection*{Определение}

Линейным пространством над полем $P$ называется множество $V$ с операциями сложения и умножения на скаляр, такими что:
\begin{enumerate}
    \item $(V,+)$ --- абелева группа;
    \item $(\alpha\beta)x=\alpha(\beta x)$;
    \item $(\alpha+\beta)x=\alpha x+\beta x$;
    \item $\alpha(x+y)=\alpha x+\alpha y$;
    \item $1x=x$.
\end{enumerate}

\subsection*{Примеры}

\begin{itemize}
    \item $P^n$;
    \item пространство многочленов $P[x]$;
    \item пространство многочленов степени не выше $n$;
    \item пространство матриц $M_n(P)$;
    \item пространство непрерывных функций $C[a,b]$ над $\R$.
\end{itemize}

\subsection*{Простейшие свойства}

\begin{enumerate}
    \item Нулевой вектор единственен.
    \item Противоположный вектор единственен.
    \item $0\cdot x=0$.
    \item $\alpha\cdot0=0$.
    \item $(-\alpha)x=\alpha(-x)=-(\alpha x)$.
\end{enumerate}

\section{Линейная зависимость и независимость}

\subsection*{Определения}

Система $a_1,\ldots,a_k$ называется \textbf{линейно зависимой}, если существуют коэффициенты $\alpha_i$, не все равные нулю, такие что
\[
    \alpha_1a_1+\cdots+\alpha_ka_k=0.
\]

Система называется \textbf{линейно независимой}, если из равенства
\[
    \alpha_1a_1+\cdots+\alpha_ka_k=0
\]
следует
\[
    \alpha_1=\cdots=\alpha_k=0.
\]

\subsection*{Критерии}

\begin{enumerate}
    \item Система линейно зависима тогда и только тогда, когда один из ее векторов является линейной комбинацией остальных.
    \item Если система содержит линейно зависимую подсистему, то вся система линейно зависима.
    \item Любая подсистема линейно независимой системы линейно независима.
    \item Если $a_1,\ldots,a_k$ линейно независимы, а $a_1,\ldots,a_k,b$ линейно зависимы, то $b\in\Lin{a_1,\ldots,a_k}$.
    \item Если $a_1,\ldots,a_k$ линейно независимы и $b\notin\Lin{a_1,\ldots,a_k}$, то $a_1,\ldots,a_k,b$ линейно независимы.
\end{enumerate}

\section{Теорема Штейница о замене}

\subsection*{Формулировка}

Пусть даны системы
\[
    a_1,\ldots,a_k
\]
и
\[
    b_1,\ldots,b_s.
\]
Если система $a_1,\ldots,a_k$ линейно независима и линейно выражается через $b_1,\ldots,b_s$, то:
\begin{enumerate}
    \item $k\le s$;
    \item в системе $b_1,\ldots,b_s$ можно выбрать $k$ векторов и заменить их на $a_1,\ldots,a_k$ так, что новая система будет эквивалентна исходной системе $b_1,\ldots,b_s$.
\end{enumerate}

\subsection*{Идея доказательства}

Так как каждый $a_i$ выражается через $b_1,\ldots,b_s$, то равенство
\[
    x_1a_1+\cdots+x_ka_k=0
\]
превращается в однородную систему из $s$ уравнений с $k$ неизвестными. Если бы $k>s$, система имела бы ненулевое решение, что противоречит линейной независимости $a_1,\ldots,a_k$.

\subsection*{Следствие}

Любые два базиса конечномерного пространства содержат одинаковое число векторов.

\section{Базисы и размерность}

\subsection*{Базис}

Система $a_1,\ldots,a_k$ называется \textbf{базисом} пространства $V$, если:
\begin{enumerate}
    \item она линейно независима;
    \item она порождает $V$:
    \[
        \Lin{a_1,\ldots,a_k}=V.
    \]
\end{enumerate}

\subsection*{Размерность}

Число векторов в любом базисе конечномерного пространства называется \textbf{размерностью}:
\[
    \dim V.
\]

Корректность определения следует из теоремы Штейница: все базисы содержат одинаковое число векторов.

\subsection*{Теорема о достройке базиса}

Любую линейно независимую систему в конечномерном пространстве можно дополнить до базиса этого пространства.

\subsection*{Идея доказательства}

Если система уже порождает пространство, она базис. Если нет, добавляем вектор, который не выражается через нее. По критерию линейной независимости система остается линейно независимой. Так как пространство конечномерно, процесс закончится.

\section{Примеры базисов}

\begin{enumerate}
    \item В $P^n$ стандартный базис:
    \[
        e_1=(1,0,\ldots,0),\ldots,e_n=(0,\ldots,0,1).
    \]
    \item В $P_n[x]$ --- пространстве многочленов степени не выше $n$:
    \[
        1,x,x^2,\ldots,x^n.
    \]
    Поэтому $\dim P_n[x]=n+1$.
    \item В $M_n(P)$ базис образуют матрицы $E_{ij}$, у которых на месте $(i,j)$ стоит $1$, а остальные элементы равны $0$. Поэтому
    \[
        \dim M_n(P)=n^2.
    \]
    \item В пространстве
    \[
        V=\{x\in\R^n\mid x_1+\cdots+x_n=0\}
    \]
    базис:
    \[
        (1,-1,0,\ldots,0),\ (0,1,-1,\ldots,0),\ldots,\ (0,\ldots,0,1,-1).
    \]
    Поэтому $\dim V=n-1$.
\end{enumerate}

\section{Выражение векторов через базис. Координаты}

Пусть $a_1,\ldots,a_n$ --- базис $V$. Тогда любой вектор $x\in V$ единственным образом представляется в виде
\[
    x=\alpha_1a_1+\cdots+\alpha_na_n.
\]

Коэффициенты $\alpha_1,\ldots,\alpha_n$ называются \textbf{координатами} вектора $x$ в базисе $a_1,\ldots,a_n$.

\subsection*{Доказательство единственности}

Пусть
\[
    x=\alpha_1a_1+\cdots+\alpha_na_n
\]
и
\[
    x=\beta_1a_1+\cdots+\beta_na_n.
\]
Тогда
\[
    (\alpha_1-\beta_1)a_1+\cdots+(\alpha_n-\beta_n)a_n=0.
\]
Так как базис линейно независим, то
\[
    \alpha_i-\beta_i=0,\qquad i=1,\ldots,n.
\]
Значит координаты единственны.

\subsection*{Свойства координат}

При сложении векторов координаты складываются, а при умножении вектора на скаляр все координаты умножаются на этот скаляр.

\section{Теорема о продолжении базиса}

Формулировка близка к теореме о достройке.

\subsection*{Теорема}

Пусть $A$ --- подпространство конечномерного пространства $V$. Тогда любой базис подпространства $A$ можно дополнить до базиса всего пространства $V$.

\subsection*{Следствие}

Для любого подпространства $A\le V$ существует подпространство $B\le V$, такое что
\[
    V=A\oplus B.
\]

\subsection*{Доказательство следствия}

Пусть $a_1,\ldots,a_k$ --- базис $A$. Дополним его до базиса $V$:
\[
    a_1,\ldots,a_k,a_{k+1},\ldots,a_n.
\]
Положим
\[
    B=\Lin{a_{k+1},\ldots,a_n}.
\]
Тогда $A+B=V$, а $A\cap B=0$ по линейной независимости общего базиса. Следовательно,
\[
    V=A\oplus B.
\]

\section{Изоморфизм линейных пространств}

\subsection*{Определение}

Линейные пространства $V$ и $W$ над полем $P$ называются \textbf{изоморфными}, если существует биекция
\[
    f:V\to W,
\]
такая что
\[
    f(\alpha x+\beta y)=\alpha f(x)+\beta f(y).
\]

\subsection*{Теорема об изоморфизме пространству строк}

Если $\dim V=n$ и $e_1,\ldots,e_n$ --- базис $V$, то отображение
\[
    \varphi:V\to P^n,
\]
которое каждому вектору ставит в соответствие столбец или строку его координат в этом базисе, является изоморфизмом.

\subsection*{Идея доказательства}

Каждый вектор имеет единственный набор координат, поэтому $\varphi$ биективно. Линейность следует из свойств координат:
\[
    [\alpha x+\beta y]=\alpha[x]+\beta[y].
\]

\section{Теорема об изоморфизме линейных пространств}

\subsection*{Теорема}

Конечномерные пространства $V$ и $W$ над одним полем $P$ изоморфны тогда и только тогда, когда
\[
    \dim V=\dim W.
\]

\subsection*{Доказательство}

Необходимость: изоморфизм переводит базис в базис, поэтому размеры базисов равны.

Достаточность: пусть
\[
    e_1,\ldots,e_n
\]
--- базис $V$, а
\[
    f_1,\ldots,f_n
\]
--- базис $W$. Зададим отображение:
\[
    T(x_1e_1+\cdots+x_ne_n)=x_1f_1+\cdots+x_nf_n.
\]
Оно линейно, биективно и поэтому является изоморфизмом.

\section{Подпространства и линейные оболочки}

\subsection*{Определение}

Непустое подмножество $U\subseteq V$ называется \textbf{подпространством}, если оно само является линейным пространством относительно операций из $V$.

\subsection*{Критерий подпространства}

Непустое множество $U\subseteq V$ является подпространством тогда и только тогда, когда
\[
    \forall x,y\in U,\ \forall \alpha,\beta\in P:\quad \alpha x+\beta y\in U.
\]

\subsection*{Примеры}

\begin{itemize}
    \item $\{0\}$ и $V$;
    \item множество решений однородной системы линейных уравнений;
    \item пространство многочленов степени не выше $n$ внутри $P[x]$.
\end{itemize}

\subsection*{Линейная оболочка}

Линейной оболочкой системы $v_1,\ldots,v_k$ называется множество всех ее линейных комбинаций:
\[
    \Lin{v_1,\ldots,v_k}
    =
    \{\alpha_1v_1+\cdots+\alpha_kv_k\mid \alpha_i\in P\}.
\]
Это наименьшее подпространство, содержащее все векторы $v_1,\ldots,v_k$.

\section{Пересечение и сумма подпространств}

\subsection*{Пересечение}

Если $U,W\le V$, то
\[
    U\cap W
\]
тоже является подпространством.

\subsection*{Сумма}

Суммой подпространств $U$ и $W$ называется
\[
    U+W=\{u+w\mid u\in U,\ w\in W\}.
\]
Она является подпространством и совпадает с линейной оболочкой $U\cup W$.

\subsection*{Прямая сумма}

Сумма называется прямой, если
\[
    U\cap W=0.
\]
Обозначение:
\[
    V=U\oplus W.
\]
Эквивалентно: каждый вектор $v\in U+W$ единственным образом представляется как
\[
    v=u+w.
\]

\subsection*{Теорема о размерности суммы}

Если $U,W$ --- конечномерные подпространства, то
\[
    \dim(U+W)=\dim U+\dim W-\dim(U\cap W).
\]

\subsection*{Идея доказательства}

Берем базис $U\cap W$ и дополняем его до базиса $U$ и до базиса $W$. Объединение полученных векторов образует базис $U+W$. Подсчет числа векторов дает формулу.

\section{Операторы линейных пространств}

Пусть $V$ и $W$ --- линейные пространства над полем $P$. Отображение $L:V\to W$ называется \textbf{линейным оператором}, если
\[
    L(\alpha x+\beta y)=\alpha L(x)+\beta L(y).
\]

\subsection*{Примеры}

\begin{itemize}
    \item тождественный оператор $I(x)=x$;
    \item нулевой оператор $O(x)=0$;
    \item оператор подобия $(\alpha I)(x)=\alpha x$;
    \item поворот плоскости;
    \item проекция $L:\R^3\to\R^2$, $L(x,y,z)=(x,y)$;
    \item дифференцирование многочленов.
\end{itemize}

\subsection*{Теорема о задании оператора на базисе}

Пусть $a_1,\ldots,a_n$ --- базис $V$, а $b_1,\ldots,b_n$ --- произвольные векторы из $W$. Тогда существует единственный линейный оператор $L:V\to W$, такой что
\[
    L(a_i)=b_i.
\]
Для $x=x_1a_1+\cdots+x_na_n$ он задается формулой
\[
    L(x)=x_1b_1+\cdots+x_nb_n.
\]

\section{Связь между базисами и координатами}

Пусть $e_1,\ldots,e_n$ --- старый базис, $f_1,\ldots,f_n$ --- новый базис. Если
\[
    f_j=c_{1j}e_1+\cdots+c_{nj}e_n,
\]
то матрица $C=(c_{ij})$ называется матрицей перехода от старого базиса к новому:
\[
    (f_1,\ldots,f_n)=(e_1,\ldots,e_n)C.
\]

Если $X$ --- старые координаты вектора, а $X'$ --- новые, то
\[
    X=CX',\qquad X'=C^{-1}X.
\]

\section{Матрица оператора}

Пусть $L:V\to W$, $a_1,\ldots,a_n$ --- базис $V$, $b_1,\ldots,b_m$ --- базис $W$.

Разложим:
\[
    L(a_i)=\alpha_{1i}b_1+\cdots+\alpha_{mi}b_m.
\]
Матрица
\[
    A_L=(\alpha_{ji})
\]
называется матрицей оператора. В $i$-м столбце стоят координаты $L(a_i)$.

Если $X$ --- координаты $x$, а $Y$ --- координаты $L(x)$, то
\[
    Y=AX.
\]

\section{Матрицы оператора в разных базисах}

Пусть $L\in\End V$. В старом базисе матрица оператора равна $A$, в новом --- $B$. Если $C$ --- матрица перехода от старого базиса к новому, то
\[
    B=C^{-1}AC.
\]

\subsection*{Подобные матрицы}

Матрицы $A$ и $B$ называются подобными, если существует обратимая $C$, такая что
\[
    B=C^{-1}AC.
\]
Подобные матрицы задают один оператор в разных базисах.

\section{Образ и ядро оператора}

Пусть $L:V\to W$.

\[
    \Image L=\{y\in W\mid \exists x\in V:\ y=L(x)\},
\]
\[
    \Ker L=\{x\in V\mid L(x)=0\}.
\]

$\Image L$ --- подпространство в $W$, $\Ker L$ --- подпространство в $V$.

\[
    r(L)=\dim\Image L,\qquad d(L)=\dim\Ker L.
\]

\subsection*{Теорема}

Если $\dim V=n$, то
\[
    r(L)+d(L)=n.
\]

Доказательство сводится к однородной системе $AX=0$: размерность пространства решений равна $n-\rank A$.

\section{Сумма и произведение операторов}

Пусть $\End_PV$ --- множество всех операторов $V\to V$.

\[
    (L_1+L_2)(x)=L_1(x)+L_2(x),
\]
\[
    (\alpha L)(x)=\alpha L(x),
\]
\[
    (L_2L_1)(x)=L_2(L_1(x)).
\]

Относительно сложения и композиции $\End_PV$ является кольцом.

Если $\dim V=n$ и базис фиксирован, то
\[
    \varphi:\End_PV\to M_n(P),\qquad \varphi(L)=A_L
\]
является изоморфизмом колец:
\[
    A_{L_1+L_2}=A_{L_1}+A_{L_2},\qquad A_{L_2L_1}=A_{L_2}A_{L_1}.
\]

\section{Алгебра над полем}

Алгебра над полем $P$ --- это множество $A$, которое является кольцом и линейным пространством над $P$, причем
\[
    \alpha(xy)=(\alpha x)y=x(\alpha y).
\]

Примеры: $M_n(P)$, $\End_PV$, $P[x]$.

Если $\dim V=n$, то
\[
    \End_PV\cong M_n(P)
\]
как алгебры. Поэтому
\[
    \dim\End_PV=n^2.
\]

\section{Мультипликативная группа кольца. Обратимые операторы}

Мультипликативная группа кольца $R$ --- множество всех обратимых элементов кольца. Обозначение: $U(R)$.

Оператор $L$ обратим, если существует $L^{-1}$, такой что
\[
    LL^{-1}=L^{-1}L=I.
\]

Множество обратимых операторов обозначается $\GL(V)$, множество обратимых матриц --- $\GL(n,P)$.

Так как оператор обратим тогда и только тогда, когда его матрица обратима,
\[
    \GL(V)\cong\GL(n,P).
\]

\section{Теорема об обратимых операторах}

Пусть $\dim V=n$, $L\in\End_PV$. Следующие условия эквивалентны:

\begin{enumerate}
    \item $L$ обратим;
    \item $L$ --- биекция;
    \item $\Image L=V$;
    \item $\Ker L=0$;
    \item для любого базиса $a_1,\ldots,a_n$ система $L(a_1),\ldots,L(a_n)$ является базисом;
    \item матрица $A_L$ обратима, то есть $\det A_L\ne0$.
\end{enumerate}

\section{Собственные векторы и собственные значения}

Ненулевой вектор $a$ называется собственным вектором оператора $L$, если
\[
    L(a)=\lambda a
\]
для некоторого $\lambda\in P$. Число $\lambda$ называется собственным значением.

Собственное подпространство:
\[
    V_\lambda=\{a\in V\mid L(a)=\lambda a\}=\Ker(L-\lambda I).
\]

\subsection*{Предложение}

Собственные векторы, отвечающие попарно различным собственным значениям, линейно независимы.

\section{Характеристический многочлен и спектр}

Пусть $A$ --- матрица оператора $L$.

Характеристическая матрица:
\[
    \lambda E-A.
\]
Характеристический многочлен:
\[
    f_L(\lambda)=\det(\lambda E-A).
\]

Спектр $\sigma(L)$ --- множество характеристических корней с учетом кратностей.

\subsection*{Независимость от базиса}

Если $B=C^{-1}AC$, то
\[
    \det(\lambda E-B)=\det(C^{-1}(\lambda E-A)C)=\det(\lambda E-A).
\]

\subsection*{Теорема}

Собственные значения оператора --- это в точности корни характеристического многочлена.

\section{Алгоритмы нахождения собственных значений и векторов}

\subsection*{Собственные значения}

\begin{enumerate}
    \item Построить матрицу $A$.
    \item Составить $\lambda E-A$.
    \item Вычислить $f(\lambda)=\det(\lambda E-A)$.
    \item Найти корни $f(\lambda)$.
\end{enumerate}

\subsection*{Собственные векторы}

Для каждого $\lambda_i$ решить систему
\[
    (\lambda_iE-A)X=0.
\]
Фундаментальная система решений дает базис собственного подпространства $V_{\lambda_i}$.

\section{Диагонализируемые операторы}

Оператор $L$ называется диагонализируемым, если существует базис, в котором его матрица диагональна.

\subsection*{Критерий}

$L$ диагонализируем тогда и только тогда, когда в $V$ существует базис из собственных векторов $L$.

Если базис $e_i$ состоит из собственных векторов и $L(e_i)=\lambda_i e_i$, то матрица в этом базисе диагональна:
\[
    \begin{pmatrix}
        \lambda_1 & & 0\\
        & \ddots & \\
        0 & & \lambda_n
    \end{pmatrix}.
\]

\section{Критерий диагонализируемости через кратности}

Пусть
\[
    f_L(\lambda)=\prod_i(\lambda-\lambda_i)^{k_i},
\]
где $k_i$ --- алгебраическая кратность собственного значения $\lambda_i$.

Всегда
\[
    \dim V_{\lambda_i}\le k_i.
\]

\subsection*{Критерий}

Оператор диагонализируем тогда и только тогда, когда для каждого $i$
\[
    \dim V_{\lambda_i}=k_i.
\]

Практически: находим корни характеристического многочлена, их кратности и размерности ядер $\Ker(\lambda_iE-A)$.

\section{Инвариантные подпространства}

Подпространство $W\le V$ называется инвариантным относительно $L$, если
\[
    L(W)\le W.
\]

\subsection*{Примеры}

\begin{itemize}
    \item $\Ker L$;
    \item $\Image L$;
    \item собственное подпространство $V_\lambda$;
    \item любое подпространство, содержащее $\Image L$.
\end{itemize}

Одномерное подпространство $\Lin{x}$ инвариантно тогда и только тогда, когда $x$ --- собственный вектор.

\subsection*{Блочные матрицы}

Если $W$ инвариантно и $\dim W=k$, то в базисе, начинающемся с базиса $W$, матрица $L$ имеет вид
\[
    \begin{pmatrix}
        A_{11} & A_{12}\\
        0 & A_{22}
    \end{pmatrix}.
\]

Если $V=W\oplus U$, где оба подпространства инвариантны, то матрица имеет блочно-диагональный вид:
\[
    \begin{pmatrix}
        A_{11} & 0\\
        0 & A_{22}
    \end{pmatrix}.
\]

\section{Индуцированный оператор}

Пусть $W\le V$ инвариантно относительно $L$. Тогда индуцированный оператор
\[
    L_W:W\to W
\]
задается формулой
\[
    L_W(x)=L(x).
\]

\subsection*{Теорема}

Характеристический многочлен индуцированного оператора $L_W$ делит характеристический многочлен оператора $L$.

\subsection*{Доказательство}

В базисе, начинающемся с базиса $W$, матрица $L$ имеет вид
\[
    A=
    \begin{pmatrix}
        A_{11} & A_{12}\\
        0 & A_{22}
    \end{pmatrix},
\]
где $A_{11}$ --- матрица $L_W$. Тогда
\[
    \det(\lambda E-A)=
    \det(\lambda E-A_{11})\det(\lambda E-A_{22}).
\]
Первый множитель --- характеристический многочлен $L_W$, значит он делит характеристический многочлен $L$.

\end{document}
